\documentclass[10pt,a4paper]{article}
\usepackage{fullpage}
\usepackage[T1]{fontenc}
\usepackage{textcomp}
\usepackage[utf8]{inputenc}
\usepackage{amsmath}
\usepackage{amssymb}
\usepackage{graphicx}
\usepackage[english]{babel}
\usepackage{url}
\usepackage{hyperref}
\usepackage{listings}
\usepackage{textcomp}
\usepackage{accsupp}
\usepackage{attachfile}
\usepackage{verbatim}
\usepackage[misc]{ifsym}
\usepackage[super]{nth}
\usepackage{dirtree}
\usepackage[os=win]{menukeys}
\usepackage{pmboxdraw}

% color def
\usepackage{color}
\definecolor{darkred}{rgb}{0.6,0.0,0.0}
\definecolor{darkgreen}{rgb}{0,0.50,0}
\definecolor{lightblue}{rgb}{0.0,0.42,0.91}
\definecolor{orange}{rgb}{0.99,0.48,0.13}
\definecolor{grass}{rgb}{0.18,0.80,0.18}
\definecolor{pink}{rgb}{0.97,0.15,0.45}

\usepackage{tikz}
\usetikzlibrary{shapes.multipart}
\usetikzlibrary {shapes.misc} 
\usetikzlibrary{shadows.blur}
\usetikzlibrary{positioning}
\usetikzlibrary{arrows}
\usetikzlibrary{matrix}
\tikzset{
stack_item/.style={rounded rectangle, rounded rectangle arc length=90,draw, anchor=center, fill=white, minimum width=1cm},
bar/.style={rounded rectangle, minimum width=1.5cm, fill=black, blur shadow={shadow blur steps=5,shadow blur extra rounding=1.3pt}},
message/.style={rectangle, fill=white, color=white, text=black},
bowling_pin/.style={circle, draw, line width=0.6mm, color=red, text=black, minimum width=8mm},
empty/.style={circle, draw, line width=0.6mm, color=white, text=black, minimum width=8mm},
table_cell/.style={rectangle, draw, color=black, text=black, minimum size=12mm}
}



% listings
\usepackage{listings}

% General Setting of listings
\lstset{
	language=Python,
	tabsize=4,
	aboveskip=1em,
	breaklines=true,
	abovecaptionskip=0pt,
	captionpos=b,
	escapeinside={\%*}{*)},
	frame=single,
	numbers=left,
	numbersep=8pt,
	numberstyle=\tiny,
	morekeywords=[1]{,as,assert,nonlocal,with,yield,self,True,False,None,} % Python builtin
	morekeywords=[2]{,__init__,__add__,__mul__,__div__,__sub__,__call__,__getitem__,__setitem__,__eq__,__ne__,__nonzero__,__rmul__,__radd__,__repr__,__str__,__get__,__truediv__,__pow__,__name__,__future__,__all__,}, % magic methods
	morekeywords=[3]{,object,type,isinstance,copy,deepcopy,zip,enumerate,reversed,list,set,len,dict,tuple,range,xrange,execfile,real,imag,reduce,str,repr,}, % common functions
	morekeywords=[4]{,Exception,NameError,IndexError,SyntaxError,TypeError,ValueError,OverflowError,ZeroDivisionError,}, % errors
	morekeywords=[5]{,ode,fsolve,sqrt,exp,sin,cos,arctan,arctan2,arccos,pi, array,norm,solve,dot,arange,isscalar,max,sum,flatten,shape,reshape,find,any,all,abs,plot,linspace,legend,quad,polyval,polyfit,hstack,concatenate,vstack,column_stack,empty,zeros,ones,rand,vander,grid,pcolor,eig,eigs,eigvals,svd,qr,tan,det,logspace,roll,min,mean,cumsum,cumprod,diff,vectorize,lstsq,cla,eye,xlabel,ylabel,squeeze,}, % numpy / math
	deletekeywords=[2]{next,set},
	deletekeywords=[3]{next,set},
	basicstyle=\ttfamily,
	backgroundcolor=\color{white},
	commentstyle=\color{darkgreen},
	keywordstyle=\color{blue}\bfseries\itshape,
	keywordstyle=[2]\color{blue}\bfseries,
	keywordstyle=[3]\color{grass},
	keywordstyle=[4]\color{red},
	keywordstyle=[5]\color{orange},
	stringstyle=\color{darkred},
	emphstyle=\color{pink}\underbar,
	morestring=[d]{"},
	showstringspaces=false,
	upquote=true,
    literate={á}{{\'a}}1
            	{é}{{\'e}}1
            	{í}{{\'{\i}}}1
            	{ó}{{\'o}}1
            	{ú}{{\'u}}1
            	{Á}{{\'A}}1
            	{É}{{\'E}}1
            	{Í}{{\'I}}1
            	{Ó}{{\'O}}1
            	{Ú}{{\'U}}1
            	{ü}{{\"u}}1
            	{Ü}{{\"U}}1
            	{ñ}{{\~n}}1
            	{Ñ}{{\~N}}1
            	{¿}{{?`}}1
            	{¡}{{!`}}1
                {█}{{\textblock}}1
                {░}{{\textltshade}}1,
}



\makeatletter
\lst@RequireAspects{writefile}

% Use \attachfile command to add listing content as paperclip.
\lstnewenvironment{attachedlisting}[2][]{%
    \lstset{#1}
	\lst@BeginAlsoWriteFile{#2.txt}%
}
{%
	\lst@EndWriteFile
	\marginpar{\attachfile[appearance=false,icon=Paperclip,mimetype=text/tex]{#2.txt}}
}

% Use \handlisting command to add file content as paperclip.
\lstnewenvironment{handlisting}[1][]{%
}
{%
	\marginpar{\attachfile[appearance=false,icon=Paperclip,mimetype=text/tex]{#1.txt}}
}

% Use accsupp package to add listing content as copyable text.
\lstnewenvironment{copyablelisting}{%
	\lst@BeginAlsoWriteFile{\jobname.txt}%
}
{%
	\lst@EndWriteFile
	\let\verbatim@processline\add@lstline
	\global\let\lstfile\empty
	\verbatiminput{\jobname.txt}%
	\marginpar{(\BeginAccSupp{method=escape,ActualText={\lstfile}}\PaperPortrait\EndAccSupp{})}
}
\def\add@lstline
{\xdef\lstfile{\unexpanded\expandafter{\lstfile}\the\verbatim@line\string^^J}}
\makeatother

\usepackage{titling}
\setlength{\droptitle}{-8em}   % This is your set screw

\title{{\huge El Juego de la Vida de Conway}}
\date{}

\begin{document}
\maketitle

\section{Introducción}

El Juego de la Vida\footnote{\href{https://es.wikipedia.org/wiki/Juego\_de\_la\_vida}{https://es.wikipedia.org/wiki/Juego\_de\_la\_vida}} es un autómata celular creado por el matemático británico John Horton Conway. Es un juego sin jugadores, lo que significa que la evolución del juego está definida por el estado inicial y no requiere más entradas. La forma de jugar es definir una configuración inicial y observar cómo evoluciona.

\section{Reglas del Juego}
El tablero de juego está definido por una malla de celdas cuadradas que se extiende hacia el infinito. Las celdas estarán en uno de dos estados, vivas o muertas. En nuestro caso, definiremos esta malla como una matriz, y los estados de las celdas como 1 para vivas y 0 para muertas. El estado de estas celdas cambiará con el tiempo, podemos llamarlo un turno o iteración. Para obtener la malla resultante para una nueva iteración, se toma la actual y se analizan los valores de cada celda y sus vecinas. Todas las celdas se actualizan al mismo tiempo en cada turno.

Las reglas son las siguientes:

\begin{description}
\item[Nace:] Si una celda muerta está rodeada exactamente por 3 celdas vivas, esa celda nacerá (estará viva en la siguiente iteración).
\item[Muere:] Una celda viva puede morir por dos razones:
    \begin{description}
    \item[Superpoblación:] Si hay más de 3 celdas vecinas vivas.
    \item[Aislamiento:] Si no hay ninguna o solo una celda vecina viva.
    \end{description} 
\item[Vive:] Una celda seguirá viva si 2 o 3 de sus vecinas están vivas.
\end{description}

El Listing~\ref{lst:example} muestra una matriz que define un tablero para este juego. Los valores 0 son celdas muertas y los 1 son celdas vivas.

\begin{attachedlisting}[caption={Ejemplo de tablero de juego.}, label=lst:example][matrix.py]
glider = [[0, 0, 0, 0, 0, 0, 0, 0, 0, 0, 0],
          [0, 0, 0, 0, 1, 0, 0, 0, 0, 0, 0],
          [0, 0, 1, 0, 1, 0, 0, 0, 0, 0, 0],
          [0, 0, 0, 1, 1, 0, 0, 0, 0, 0, 0],
          [0, 0, 0, 0, 0, 0, 0, 0, 0, 0, 0],
          [0, 0, 0, 0, 0, 0, 0, 0, 0, 0, 0],
          [0, 0, 0, 0, 0, 0, 0, 0, 0, 0, 0],
          [0, 0, 0, 0, 0, 0, 0, 0, 0, 0, 0],
          [0, 0, 0, 0, 0, 0, 0, 0, 0, 0, 0],
          [0, 0, 0, 0, 0, 0, 0, 0, 0, 0, 0],
          [0, 0, 0, 0, 0, 0, 0, 0, 0, 0, 0]]
\end{attachedlisting}

\newpage
\section{Creando el Programa}

Crea las funciones descritas para completar un programa que simule el juego de la vida.

\subsection{\texttt{read\_board(path)}}

Implementa una función \lstinline|read_board(path: str) -> list[list[int]]|. Esta función tomará una ruta de archivo como argumento y devolverá la matriz que contiene. Tendrás que verificar si el archivo realmente existe y si la matriz que contiene es válida, es decir, si todas las filas contienen la misma cantidad de columnas.

\subsection{\texttt{print\_matrix(matrix)}}

Implementa una función \lstinline|print_matrix(matrix: list[list[int]]) -> None| que tome un tablero de juego como parámetro y lo imprima en la consola. Para imprimirlo, usaremos dos bloques oscuros, █, para las celdas que están vivas y dos bloques sombreados, ░, para las celdas que están muertas. Se necesitan dos para hacer que cada celda sea cuadrada.

Si el tablero es demasiado grande, es posible que tengas que desanclar la consola.

La Sección~\ref{sec:improve} tiene una alternativa sobre cómo mostrar la matriz que podría generar una pantalla de mejor apariencia.

Ejemplo de impresión de una celda viva y una muerta. \textbf{En el \texttt{pdf} no se puede ver correctamente, consulta el archivo adjunto}.
\begin{attachedlisting}[caption={Sentencias de impresión para los caracteres sombreados.}][print.txt]
print("██") # Viva
print("░░") # Muerta
\end{attachedlisting}

La salida esperada de esta función con la matriz anterior se verá de la siguiente manera:

\begin{lstlisting}
print_matrix(glider)
# Salida: 
# ░░░░░░░░░░░░░░░░░░░░░░
# ░░░░░░░░██░░░░░░░░░░░░
# ░░░░██░░██░░░░░░░░░░░░
# ░░░░░░████░░░░░░░░░░░░
# ░░░░░░░░░░░░░░░░░░░░░░
# ░░░░░░░░░░░░░░░░░░░░░░
# ░░░░░░░░░░░░░░░░░░░░░░
# ░░░░░░░░░░░░░░░░░░░░░░
# ░░░░░░░░░░░░░░░░░░░░░░
# ░░░░░░░░░░░░░░░░░░░░░░
# ░░░░░░░░░░░░░░░░░░░░░░
\end{lstlisting}

\subsection{\texttt{get\_neighbours(matrix, row, column)}}

Implementa una función \lstinline|get_neighbours(matrix: list[list[int]], row: int, column: int) -> list[int]|. Esta función tomará una matriz y, como parámetros, índices de fila y columna. Devolverá una lista que contiene los valores de los vecinos de la celda representada por los índices de fila y columna. Recuerda que cada celda en la matriz tendrá 8 vecinos. Para simplificar esta función, si los valores de fila/columna son los primeros o los últimos, la función devolverá una lista vacía (ignoraremos el borde de la matriz).

\begin{lstlisting}
print(get_neighbours(glider, 2, 2))
# Salida: [0, 0, 0, 0, 0, 0, 0, 1]
print(get_neighbours(glider, 0, 1))
# Salida: []
\end{lstlisting}

\subsection{\texttt{count\_alive(neighbours)}}

Implementa una función \lstinline|count_alive(neighbours: list[int]) -> int| que, tomando una lista de vecinos como parámetros, devuelva cuántos de esos vecinos están vivos.

\begin{lstlisting}
print(count_alive(get_neighbours(glider, 2,2)))
# Salida: 1
\end{lstlisting}

\subsection{\texttt{get\_next\_iteration(matrix)}}

Implementa una función \lstinline|get_next_iteration(matrix: list[list[int]]) -> list[list[int]]| que tome una matriz de tablero como parámetro y devuelva el tablero con la siguiente iteración según las reglas del juego de la vida. Crea una matriz llena de 0s y luego verifica para cada elemento en la matriz cuántos de sus vecinos están vivos. Establece los valores de la nueva matriz según las siguientes reglas:
\begin{itemize}
\item Si la celda está viva
    \begin{itemize}
    \item La celda seguirá viva si hay 2 o 3 vecinos vivos (valor a 1)
    \item Si hay superpoblación (más de 3) o subpoblación (menos de 2), la celda muere  
    \end{itemize}
\item si la celda está muerta
    \begin{itemize}
    \item Si hay exactamente 3 vecinos vivos, la celda nacerá
    \end{itemize}
\end{itemize}

La función debería funcionar de la siguiente manera:

\begin{lstlisting}
print_matrix(glider)
# Salida: 
# ░░░░░░░░░░░░░░░░░░░░░░
# ░░░░░░░░██░░░░░░░░░░░░
# ░░░░██░░██░░░░░░░░░░░░
# ░░░░░░████░░░░░░░░░░░░
# ░░░░░░░░░░░░░░░░░░░░░░
# ░░░░░░░░░░░░░░░░░░░░░░
# ░░░░░░░░░░░░░░░░░░░░░░
# ░░░░░░░░░░░░░░░░░░░░░░
# ░░░░░░░░░░░░░░░░░░░░░░
# ░░░░░░░░░░░░░░░░░░░░░░
# ░░░░░░░░░░░░░░░░░░░░░░
new = get_next_iteration(glider)
print_matrix(new)
# Salida:
# ░░░░░░░░░░░░░░░░░░░░░░
# ░░░░░░██░░░░░░░░░░░░░░
# ░░░░░░░░████░░░░░░░░░░
# ░░░░░░████░░░░░░░░░░░░
# ░░░░░░░░░░░░░░░░░░░░░░
# ░░░░░░░░░░░░░░░░░░░░░░
# ░░░░░░░░░░░░░░░░░░░░░░
# ░░░░░░░░░░░░░░░░░░░░░░
# ░░░░░░░░░░░░░░░░░░░░░░
# ░░░░░░░░░░░░░░░░░░░░░░
# ░░░░░░░░░░░░░░░░░░░░░░
\end{lstlisting}

\subsection{\texttt{count\_changes(old, new)}}

Implementa una función \lstinline|count_changes(old: list[list[int]], new: list[list[int]]) -> (int, int)|. Esta función tomará las matrices de dos iteraciones y devolverá dos enteros, el número de nacimientos y muertes que han ocurrido de una iteración a la otra.

Un ejemplo del uso de la función es el siguiente:

\begin{lstlisting}
print_matrix (glider)
print("La siguiente iteración es:")
print()
new = get_next_iteration(glider)
print_matrix(new)
births, deaths = count_changes(glider, new))
print(f"En esta iteración ha habido {deaths} muertes y {births} nacimientos")
# Salida completa:
# ░░░░░░░░░░░░░░░░░░░░░░
# ░░░░░░░░██░░░░░░░░░░░░
# ░░░░██░░██░░░░░░░░░░░░
# ░░░░░░████░░░░░░░░░░░░
# ░░░░░░░░░░░░░░░░░░░░░░
# ░░░░░░░░░░░░░░░░░░░░░░
# ░░░░░░░░░░░░░░░░░░░░░░
# ░░░░░░░░░░░░░░░░░░░░░░
# ░░░░░░░░░░░░░░░░░░░░░░
# ░░░░░░░░░░░░░░░░░░░░░░
# ░░░░░░░░░░░░░░░░░░░░░░
# La siguiente iteración es:
# 
# ░░░░░░░░░░░░░░░░░░░░░░
# ░░░░░░██░░░░░░░░░░░░░░
# ░░░░░░░░████░░░░░░░░░░
# ░░░░░░████░░░░░░░░░░░░
# ░░░░░░░░░░░░░░░░░░░░░░
# ░░░░░░░░░░░░░░░░░░░░░░
# ░░░░░░░░░░░░░░░░░░░░░░
# ░░░░░░░░░░░░░░░░░░░░░░
# ░░░░░░░░░░░░░░░░░░░░░░
# ░░░░░░░░░░░░░░░░░░░░░░
# ░░░░░░░░░░░░░░░░░░░░░░
# En esta iteración ha habido 2 muertes y 2 nacimientos
\end{lstlisting}

\subsection{Completa el código}

Comienza creando la función \lstinline|main()| para el programa. La primera tarea es definir la configuración inicial del tablero. Tienes varias opciones para esto: puedes usar el ejemplo proporcionado en este documento, leer la configuración desde un archivo o incluso diseñar tu propio tablero personalizado. Una vez configurado el tablero inicial, el programa entrará en un bucle donde calculará y mostrará continuamente la siguiente iteración del tablero.

Para asegurar que el programa se detenga en cada iteración, usa una instrucción \lstinline|input()|. Esto pedirá al usuario que decida si quiere continuar o detener el programa. Durante cada iteración, necesitarás llevar un registro de dos estados del tablero simultáneamente: la iteración actual y la iteración anterior. Esto es crucial para calcular los cambios con precisión.

En cada iteración, muestra el número de iteración junto con el número de muertes y nacimientos que ocurrieron. Además, mantén un total acumulado de todas las muertes y nacimientos durante la ejecución del programa. Cuando el usuario decida detener el programa, presentarás estos totales acumulados para proporcionar una visión completa de la evolución del tablero.

Siguiendo estos pasos, crearás una simulación dinámica que permitirá a los usuarios observar la progresión del estado del tablero a lo largo de múltiples iteraciones.

El resultado esperado es el siguiente:

\begin{lstlisting}
# Iniciando el juego
# Usando la matriz inicial:
# ░░░░░░░░░░░░░░░░░░░░░░
# ░░░░░░░░██░░░░░░░░░░░░
# ░░░░██░░██░░░░░░░░░░░░
# ░░░░░░████░░░░░░░░░░░░
# ░░░░░░░░░░░░░░░░░░░░░░
# ░░░░░░░░░░░░░░░░░░░░░░
# ░░░░░░░░░░░░░░░░░░░░░░
# ░░░░░░░░░░░░░░░░░░░░░░
# ░░░░░░░░░░░░░░░░░░░░░░
# ░░░░░░░░░░░░░░░░░░░░░░
# ░░░░░░░░░░░░░░░░░░░░░░
# 
# Matriz para la iteración 1
# ░░░░░░░░░░░░░░░░░░░░░░
# ░░░░░░██░░░░░░░░░░░░░░
# ░░░░░░░░████░░░░░░░░░░
# ░░░░░░████░░░░░░░░░░░░
# ░░░░░░░░░░░░░░░░░░░░░░
# ░░░░░░░░░░░░░░░░░░░░░░
# ░░░░░░░░░░░░░░░░░░░░░░
# ░░░░░░░░░░░░░░░░░░░░░░
# ░░░░░░░░░░░░░░░░░░░░░░
# ░░░░░░░░░░░░░░░░░░░░░░
# ░░░░░░░░░░░░░░░░░░░░░░
# 
# En esta iteración ha habido 2 muertes y 2 nacimientos
# Introduce 's' para detener, cualquier otro carácter para continuar: 
# Matriz para la iteración 2
# ░░░░░░░░░░░░░░░░░░░░░░
# ░░░░░░░░██░░░░░░░░░░░░
# ░░░░░░░░░░██░░░░░░░░░░
# ░░░░░░██████░░░░░░░░░░
# ░░░░░░░░░░░░░░░░░░░░░░
# ░░░░░░░░░░░░░░░░░░░░░░
# ░░░░░░░░░░░░░░░░░░░░░░
# ░░░░░░░░░░░░░░░░░░░░░░
# ░░░░░░░░░░░░░░░░░░░░░░
# ░░░░░░░░░░░░░░░░░░░░░░
#
# En esta iteración ha habido 2 muertes y 2 nacimientos
# Introduce 's' para detener, cualquier otro carácter para continuar: 
# Matriz para la iteración 3
# ░░░░░░░░░░░░░░░░░░░░░░
# ░░░░░░░░░░░░░░░░░░░░░░
# ░░░░░░██░░██░░░░░░░░░░
# ░░░░░░░░████░░░░░░░░░░
# ░░░░░░░░██░░░░░░░░░░░░
# ░░░░░░░░░░░░░░░░░░░░░░
# ░░░░░░░░░░░░░░░░░░░░░░
# ░░░░░░░░░░░░░░░░░░░░░░
# ░░░░░░░░░░░░░░░░░░░░░░
# ░░░░░░░░░░░░░░░░░░░░░░
# ░░░░░░░░░░░░░░░░░░░░░░
# 
# En esta iteración ha habido 2 muertes y 2 nacimientos
# Introduce 's' para detener, cualquier otro carácter para continuar: s
# En total ha habido 6 muertes y 6 nacimientos
\end{lstlisting}

\section{Mejoras adicionales}
\label{sec:improve}

\begin{itemize}
\item Usa \texttt{colorama} para mejorar la salida del tablero. Puedes usar dos espacios vacíos, en lugar de los caracteres sombreados, configurados con un color de fondo para mostrar las células que están muertas o vivas.
\item Crea un menú para el programa que le dé al usuario la opción de elegir entre un conjunto de archivos en una carpeta. Obtén los archivos que están en la carpeta usando \lstinline|os.listdir()| del módulo \lstinline|os|.
\item En lugar de usar la función \lstinline|input()| para pausar después de cada iteración, considera usar la función \lstinline|time.sleep()| del módulo \lstinline|time|. Esta función toma un número que representa la duración en segundos y pausa el programa durante ese tiempo antes de continuar. Si necesitas especificar la duración en milisegundos, simplemente divide el número por 1000.

Perderás la capacidad de detener el programa presionando ``s''. ¿Hay otra forma en la que podamos detener el programa de manera elegante? Podrías tener que verificar las pulsaciones de teclas.
\end{itemize}


\end{document}